\chapter{$D_\lambda$-模的上同调}
\label{ch:02.2}

\section{旗簇\index{qicu@旗簇|hyperpage}上的齐性扭微分算子层}
\label{sec:02.2.1}

本节把第 1 章关于齐性扭微分算子层的构造，用于「连通半单代数群 $G$ 作用在其旗簇 $X$ 上」这一情形。

设 $\g$ 是复半单李代数，$G$ 是 $\g$ 的内自同构群。则 $\g$ 的旗簇 $X$ 可与 $\g$ 的全体 Borel 子代数构成的簇等同。群 $G$ 自然作用在平凡向量丛 $X\times\g\to X$ 上；Borel 子代数构成的\emph{重言向量丛}（tautological vector bundle）$\mathcal{B}$ 是它的齐性向量子丛。对每个 $x\in X$，记 $x$ 对应的 Borel 子代数为 $\bb_x$，其幂零根\index{milinggen@幂零根|hyperpage}为 $\n_x$。于是幂零根构成 $\mathcal{B}$ 的齐性向量子丛 $\mathcal{N}$。此外，记 $B_x$ 为 $G$ 中对应于 $\bb_x$ 的 Borel 子群。则 $B_x$ 是 $x$ 在 $G$ 中的稳定化子。

置 $H=\mathcal{B}/\mathcal{N}$。则 $H$ 是 $X$ 上的齐性向量丛，在 $x\in X$ 处的纤维为 $\h_x=\bb_x/\n_x$。群 $B_x$ 在 $\bb_x/\n_x$ 上平凡地作用，故 $H$ 是 $X$ 上的平凡向量丛；其整体截面\index{zhengtijiemian@整体截面|hyperpage}构成的李代数 $\h$ 与 $\bb_x/\n_x$ 自然同构（对任意 $x\in X$ 皆然）。称这个交换李代数 $\h$ 为 $\g$ 的 Cartan 代数。

设 $\calO_X$ 是代数簇 $X$ 的结构层。仍如第 1 章 §2，置 $\g^\circ=\calO_X\otimes_\C\g$，它是平凡丛 $X\times\g$ 的局部截面层。记 $\bb^\circ$ 与 $\n^\circ$ 分别为 $\mathcal{B}$ 与 $\mathcal{N}$ 的局部截面构成的子层。设 $\tau$ 是李代数 $\g$ 到 $X$ 上向量场李代数的自然同态。在 $\g^\circ$ 上置
$$[f\otimes\xi,\,g\otimes\eta]=f\tau(\xi)g\otimes\eta-g\tau(\eta)f\otimes\xi+fg\otimes[\xi,\eta],\qquad f,g\in\calO_X,\ \xi,\eta\in\g,$$
就使 $\g^\circ$ 成为复李代数层。若把 $\tau$ 延拓为 $\g^\circ$ 到 $X$ 上局部向量场李代数层的自然同态，则 $\ker\tau$ 恰为 $\bb^\circ$。此外还有下述结果。

\begin{lemma}\label{lem:02.2.1.1}
(i) 层 $\bb^\circ$ 是 $\g^\circ$ 中的理想层，$\bb^\circ$ 上的换位子是 $\calO_X$-线性的。

(ii) 层 $\n^\circ$ 是 $\g^\circ$ 中的理想层。
\end{lemma}

\begin{proof}
(i) 的第一个断言由 $\bb^\circ=\ker\tau$ 得出。又若 $\sum_i f_i\otimes\xi_i\in\bb^\circ$，$g\in\calO_X$，$\eta\in\g$，则
\begin{align*}
\Bigl[\sum_i f_i\otimes\xi_i,\,g\otimes\eta\Bigr]
&=\sum_i f_ig\otimes[\xi_i,\eta]+\sum_i f_i\tau(\xi_i)g\otimes\eta-\sum_i g\tau(\eta)f_i\otimes\xi_i\\
&=g\Bigl[\sum_i f_i\otimes\xi_i,\,1\otimes\eta\Bigr],
\end{align*}
这立即给出第二个断言。由同一公式，为证 (ii) 只需验证：对任意 $\sum_i f_i\otimes\xi_i\in\n^\circ$ 与 $\eta\in\g$，换位子 $[\sum_i f_i\otimes\xi_i,1\otimes\eta]$ 属于 $\n^\circ$。由齐性，对任意 $g\in G$ 有 $g(\sum_i f_i\otimes\xi_i)\in\n^\circ$。对此求导，即得
$$\sum_i\tau(\eta)f_i\otimes\xi_i+\sum_i f_i\otimes[\eta,\xi_i]\in\n^\circ\qquad(\eta\in\g);$$
而按换位子的定义，这个表达式等于 $[1\otimes\eta,\sum_i f_i\otimes\xi_i]$。
\end{proof}

商层\index{shangceng@商层|hyperpage} $\h^\circ=\bb^\circ/\n^\circ$ 是 $H$ 的局部截面层，因此等于交换李代数层 $\calO_X\otimes_\C\h$。

同样，第 1 章 §2 用
$$(f\otimes\xi)(g\otimes\eta)=f\tau(\xi)g\otimes\eta+fg\otimes\xi\eta,\qquad f,g\in\calO_X,\ \xi\in\g,\ \eta\in U(\g),$$
在层 $U^\circ=\calO_X\otimes_\C U(\g)$ 上定义了乘法。于是 $U^\circ$ 成为 $X$ 上的复结合代数层。显然，$\g^\circ$ 是 $U^\circ$ 的子层，而 $U^\circ$ 中的自然换位子诱导了 $\g^\circ$ 上的换位子。由 \ref{lem:02.2.1.1} 可知，$\n^\circ$ 在 $U^\circ$ 中生成的右理想层 $\n^\circ U^\circ$ 是 $U^\circ$ 的双边理想层。因此商 $\D_\h=U^\circ/\n^\circ U^\circ$ 是 $X$ 上的复结合代数层。

$\g^\circ$ 到 $\D_\h$ 的自然态射诱导了李子代数层 $\h^\circ$ 到 $\D_\h$ 的态射，于是有 $\h$ 的包络代数 $U(\h)$ 到 $\D_\h$ 的整体截面的自然同态 $\varphi$。群 $G$ 在结构层 $\calO_X$ 与 $U(\g)$ 上的作用诱导了 $U^\circ$ 与 $\D_\h$ 上的自然 $G$-作用。另一方面，$H$ 的平凡性意味着诱导出的 $G$-作用在 $\h$ 上平凡。由此可见，$\varphi$ 把 $U(\h)$ 映入 $\Gamma(X,\D_\h)$ 的 $G$-不变元。

\begin{lemma}\label{lem:02.2.1.2}
(i) $U(\h)$ 到「$\Gamma(X,\D_\h)$ 中全部 $G$-不变元构成的子代数」的自然态射 $\varphi$ 是单射。

(ii) $\varphi$ 的像含于 $\D_\h$ 的中心。
\end{lemma}

\begin{proof}
(i) 设 $x\in X$。则 $\calO_X$-模 $\D_\h$ 在 $x$ 处的几何纤维 $T_x(\D_\h)$ 等于 $U(\g)/\n_xU(\g)$。把 $\varphi$ 与「在 $x$ 处取截面之值」复合，所得对应自然映射
$$U(\h)\longrightarrow U(\bb_x)/\n_xU(\bb_x)\longrightarrow U(\g)/\n_xU(\g),$$
由 Poincar\'e--Birkhoff--Witt 定理，它是单射。因此 $\varphi$ 是单射。

(ii) 对 $G$-作用求导可知，$\varphi(U(\h))$ 中的元素与 $\g$ 在 $\D_\h$ 中的像交换。而 $\D_\h$ 由 $\calO_X$ 与 $\g$ 的像生成，断言由此得出。
\end{proof}

设 $x\in X$，$\bb_x$ 是 $x$ 对应的 Borel 子代数。设 $\bar\n$ 是与 $\bb_x$ 相对的某个 Borel 子代数的幂零根，$\bar N$ 是 $G$ 中对应的连通子群。则由 Bruhat 分解 [2, 14.11]，轨道映射
$$\bar N\longrightarrow X,\qquad \bar n\longmapsto\bar n x$$
是簇 $\bar N$ 到 $x$ 的某个开邻域 $U\subset X$ 的同构。设 $s:U\to\bar N$ 是其逆映射。显然，$U(\bar\n)$ 到 $U(\g)$ 的包含诱导了代数层 $\calO_U\otimes_\C U(\bar\n)$ 到 $U^\circ|_U$ 的单态射\index{dantaishe@单态射|hyperpage}。于是有代数层 $\calO_U\otimes_\C U(\bar\n)$ 到 $\D_\h$ 的自然态射。

再者，若把张量积 $(\calO_U\otimes_\C U(\bar\n))\otimes_\C U(\h)$ 看作代数层，则由上述讨论，有从 $(\calO_U\otimes_\C U(\bar\n))\otimes_\C U(\h)$ 到 $\D_\h|_U$ 的代数层自然态射 $\psi$。

\begin{lemma}\label{lem:02.2.1.3}
态射
$$\psi:(\calO_U\otimes_\C U(\bar\n))\otimes_\C U(\h)\longrightarrow\D_\h|_U$$
是代数层的同构。
\end{lemma}

\begin{proof}
与 \ref{lem:02.2.1.2}(i) 的证明一样，可得：把 $\psi$ 与「在 $u\in U$ 处取值」的映射复合，所得的态射对应于「$(\calO_U\otimes_\C U(\bar\n))\otimes_\C U(\h)$ 在 $u$ 处的求值映射」与自然线性映射 $U(\bar\n)\otimes_\C U(\h)\to U(\g)/\n_uU(\g)$ 的复合。由 Poincar\'e--Birkhoff--Witt 定理，后一映射是单射。由此 $\psi$ 是单射。

还需证 $\psi$ 是满射。显然，$\D_\h|_U$ 作为代数层由 $\calO_U$ 与 $\g$ 在 $\D_\h|_U$ 中的像生成。另一方面，作为向量空间，对任意 $u\in U$ 有 $\g=\bar\n\oplus\bb_u$，于是有良定义的线性同构 $\zeta(u):\g\to\g$：它在 $\bar\n$ 上是恒等，在 $\bb_x$ 上是 $\operatorname{Ad}s(u)$。因此每个 $\xi\in\g$ 决定 $\g^\circ$ 在 $U$ 上的截面 $\zeta_\xi:u\mapsto\zeta(u)\xi$。由此可见，$\D_\h|_U$ 由 $\calO_U$ 以及截面 $\zeta_\xi$（$\xi\in\g$）在 $\D_\h|_U$ 中的像生成。但若 $\xi\in\bar\n$，则对任意 $u\in U$ 有 $\zeta(u)\xi=\xi$，故该截面在 $\psi$ 的像中；若 $\xi\in\bb_x$，相应的截面也在 $\psi$ 的像中。于是 $\psi$ 也是满射。
\end{proof}

特别地，若把 $\D_\h$ 看作 $U(\h)$-模，则有下述推论。

\begin{corollary}\label{cor:02.2.1.4}
$U(\h)$-模 $\D_\h$ 是局部自由的。
\end{corollary}

此外，还可以加强 \ref{lem:02.2.1.2}(i)。

\begin{lemma}\label{lem:02.2.1.5}
$U(\h)$ 到「$\Gamma(X,\D_\h)$ 中全部 $G$-不变元构成的子代数」的自然态射 $\varphi$ 是同构。
\end{lemma}

\begin{proof}
由 \ref{lem:02.2.1.2}(i) 知 $\varphi$ 是单射。设 $s$ 是 $\D_\h$ 的 $G$-不变整体截面，则 $s$ 在 $x$ 处的值必须是 $B_x$-不变的。这意味着：若取定 $\bb_x$ 中的 Cartan 子代数 $\mathfrak{c}$，则 $s(x)$ 在 $U(\g)/\n_xU(\g)$ 中关于 $\mathfrak{c}$ 的权为零。因此 $s(x)$ 在 $U(\h)$ 的像中，即存在 $t\in\varphi(U(\h))$，使 $t-s$ 是 $G$-不变截面且在 $x$ 处为零。由 $G$-不变性，$t-s$ 在 $X$ 的每一点都为零。由 \ref{lem:02.2.1.3}，$\D_\h$ 作为关于左乘的 $\calO_X$-模是局部自由的，故 $t-s=0$，即 $s=t\in\varphi(U(\h))$。
\end{proof}

另一方面，有 $U(\g)$ 到 $\D_\h$ 的自然同态，它诱导 $U(\g)$ 的中心 $\zz(\g)$ 到 $\Gamma(X,\D_\h)$ 的自然同态。这个同态的像含于 $\Gamma(X,\D_\h)$ 的 $G$-不变元子代数，故由 \ref{lem:02.2.1.5}，它含于 $\varphi(U(\h))$。最后，有典范的 Harish-Chandra 同态
$$\gamma:\zz(\g)\longrightarrow U(\h)$$
（见 [5, Ch. VII, §6, no. 4]），其定义如下。首先，对任意 $x\in X$，中心 $\zz(\g)$ 含于子代数 $U(\bb_x)$ 与 $U(\g)$ 的右理想 $\n_xU(\g)$ 之和中。因此有自然投影
$$\zz(\g)\longrightarrow U(\bb_x)\big/\bigl(\n_xU(\g)\cap U(\bb_x)\bigr)=U(\bb_x)/\n_xU(\bb_x)=U(\h_x).$$
它与 $U(\h_x)$ 到 $U(\h)$ 的自然同构的复合与 $x$ 无关，按定义就等于 $\gamma$。

\begin{proposition}\label{prop:02.2.1.6}
由自然代数同态构成的图表
\begin{center}
\begin{tikzcd}
\zz(\g) \arrow[rr,"\gamma"] \arrow[dr] & & U(\h) \arrow[dl,"\varphi"]\\
 & \Gamma(X,\D_\h) &
\end{tikzcd}
\end{center}
交换。
\end{proposition}

\begin{proof}
由 \ref{lem:02.2.1.3}，$\D_\h$ 作为关于左乘的 $\calO_X$-模是局部自由的。因此只需证明：对任意 $x\in X$，$\varphi\circ\gamma$ 与「$\zz(\g)$ 到 $\D_\h$ 的典范同态」分别与求值映射复合所得的结果相同。但这立即由 $T_x(\D_\h)=U(\g)/\n_xU(\g)$ 得出。
\end{proof}

设 $x\in X$，取定 $\bb_x$ 中的 Cartan 子代数 $\mathfrak{c}$。设 $R$ 是 $\g$ 关于 $\mathfrak{c}$ 的根系\index{genxi@根系|hyperpage}（在 $\mathfrak{c}^*$ 中），并对 $\alpha\in R$ 置
$$\g_\alpha=\{\xi\in\g\mid [\eta,\xi]=\alpha(\eta)\xi,\ \eta\in\mathfrak{c}\},$$
称之为根 $\alpha$ 决定的根子空间。选取正根\index{zhenggen@正根|hyperpage}集
$$R^+=\{\alpha\in R\mid \g_\alpha\subset\n_x\},$$
就在 $R$ 上定义了序。于是典范同构 $\mathfrak{c}\to\h_x\to\h$ 诱导三元组 $(\mathfrak{c}^*,R,R^+)$ 到三元组 $(\h^*,\Sigma,\Sigma^+)$ 的同构，其中 $\Sigma$ 是 $\h^*$ 中的根系，$\Sigma^+$ 是 $\Sigma$ 的一组正根。显然，$\Sigma$ 与 $\Sigma^+$ 与 $x\in X$ 的选取无关。称三元组 $(\h^*,\Sigma,\Sigma^+)$ 为 $\g$ 的\emph{Cartan 三元组}（Cartan triple）；并把 Cartan 三元组 $(\h^*,\Sigma,\Sigma^+)$ 到 $(\mathfrak{c}^*,R,R^+)$ 的逆同构称为在 $x$ 处的\emph{特殊化}（specialization）。

设 $W$ 是 $\Sigma$ 的 Weyl 群，$\rho$ 是 $\Sigma$ 中全部正根之和的一半。$\h$ 的包络代数 $U(\h)$ 自然同构于 $\h^*$ 上的多项式代数，故每个 $\lambda\in\h^*$ 决定 $U(\h)$ 到 $\C$ 的同态。设 $I_\lambda$ 是由 $\lambda+\rho$ 决定的同态 $\phi_\lambda:U(\h)\to\C$ 的核。
\begin{annote}
底本只把 $I_\lambda$ 说成「由 $\lambda+\rho$ 决定的同态 $\phi_\lambda$ 的核」，没有写出它的生成元。把 $U(\h)$ 与 $\h^*$ 上的多项式代数等同后，$I_\lambda$ 就是由 $h-(\lambda+\rho)(h)\cdot1$（$h\in\h$）生成的双边理想。正文随即给出的 $I_{-\rho}=\h U(\h)$ 正与此相符：$\lambda=-\rho$ 时 $\lambda+\rho=0$，生成元退化为 $h$。[来源：底本 p.20，上下文自证]
\end{annote}
则 $\gamma^{-1}(I_\lambda)$ 是 $\zz(\g)$ 中的极大理想；又由 Harish-Chandra 的一个结果 [5, Ch. VIII, §8, Cor. 1 of Th. 2]，对 $\lambda,\mu\in\h^*$，
$$\gamma^{-1}(I_\lambda)=\gamma^{-1}(I_\mu)\quad\text{当且仅当}\quad \text{存在 } w\in W \text{ 使 } w\lambda=\mu.$$

对任意 $\lambda\in\h^*$，由 \ref{lem:02.2.1.3}，层 $I_\lambda\D_\h$ 是 $\D_\h$ 中的双边理想层；因此
$$\D_\lambda=\D_\h/I_\lambda\D_\h$$
是 $X$ 上的复结合代数层。当 $\lambda=-\rho$ 时，$I_{-\rho}=\h U(\h)$，故 $\D_{-\rho}=U^\circ/\bb^\circ U^\circ$，即它是 $X$ 上的局部微分算子层。一般的 $\D_\lambda$（$\lambda\in\h^*$）都是 $X$ 上的齐性扭微分算子层。这可由……推出，也可直接由 \ref{lem:02.2.1.3} 推出。在我们于……所用的扭微分算子层的参数化中，
$$\D_\lambda=\D_{X,\lambda+\rho},\qquad \lambda\in\h^*.$$
\begin{annote}
底本在此约定 $\D_\lambda=\D_{X,\lambda+\rho}$，即第 2 章所用的参数 $\lambda$ 比第 1 章 §2 中 $\D_{X,\lambda}$ 的参数小了 $\rho$。这是作者有意的约定，也最容易看漏：此后 $\D_\lambda$ 一律指 $\D_{X,\lambda+\rho}$，而不是 $\D_{X,\lambda}$。正文所说的 $\D_{-\rho}=U^\circ/\bb^\circ U^\circ$ 正与此相符。[来源：底本 p.20，上下文自证]
\end{annote}

设 $\theta$ 是 $\h^*$ 中的 Weyl 群轨道，$\lambda\in\theta$。记 $J_\theta=\gamma^{-1}(I_\lambda)$ 为 $\theta$ 决定的 $\zz(\g)$ 中的极大理想。记 $\chi_\lambda$ 为 $\zz(\g)$ 到 $\C$ 的、核为 $J_\theta$ 的同态。如前述，$\chi_\lambda$ 只依赖于 $\lambda$ 的 Weyl 群轨道 $\theta$。$J_\theta$ 中的元素映入 $\D_\lambda$ 的零截面。因此有典范态射
$$U_\theta=U(\g)/J_\theta U(\g)\longrightarrow\Gamma(X,\D_\lambda).$$
\begin{annote}
底本只说 $J_\theta$ 是 $\zz(\g)$ 中的极大理想，没有说明 $J_\theta U(\g)$ 指什么。这里 $J_\theta U(\g)$ 是 $U(\g)$ 中由 $J_\theta$ 生成的双边理想，故 $U_\theta=U(\g)/J_\theta U(\g)$ 正是把 $U(\g)$ 对「无穷小特征标 $\chi_\lambda$ 对应的中心理想」取商所得的代数；$\theta$ 跑遍 Weyl 群轨道时，恰好把中心特征标相同、而权不同的 $\lambda$ 归为一体。[来源：底本 p.20，上下文自证]
\end{annote}
我们将在 §6 看到，这个态射实际上是同构。

$U_\theta$-模构成的范畴 $\mathcal{M}(U_\theta)$ 的对象，也可看作无穷小特征标为 $\chi_\lambda$ 的 $U(\g)$-模。全体 $\D_\lambda$-模构成的范畴 $\mathcal{M}(\D_\lambda)$ 有足够多的内射对象 [8, III.2.2]。此外，内射 $\D_\lambda$-模是\emph{松弛}（flasque）的 [8, III.2.4]。这意味着 $\D_\lambda$-模 $V$ 的上同调模 $H^i(X,V)$ 带有 $\Gamma(X,\D_\lambda)$-模的自然结构。特别地，由前述，它们可看作 $U_\theta$-模。于是得到一族函子
$$H^i(X,-):\mathcal{M}(\D_\lambda)\longrightarrow\mathcal{M}(U_\theta),\qquad 0\le i\le\dim X.$$
下面几节将研究它们的基本性质。

\section{$D_\lambda$-模的平移原理\index{pingyiyuanli@平移原理|hyperpage}}
\label{sec:02.2.2}

本节收集研究 $\D_\lambda$-模的上同调所需的若干技术结果。

设 $Q(\Sigma)$ 是 $\h^*$ 中的根格。对任意 $\lambda\in\h^*$，记 $W_\lambda$ 为 $W$ 的子群
$$W_\lambda=\{w\in W\mid w\lambda-\lambda\in Q(\Sigma)\}.$$

设 $\Sigma^\vee$ 是 $\h$ 中与 $\Sigma$ 对偶的根系；对任意 $\alpha\in\Sigma$，记 $\alpha^\vee\in\Sigma^\vee$ 为 $\alpha$ 的对偶根。则由 [5, Ch. VI, §2, Ex. 2]，$W_\lambda$ 是根系
$$\Sigma_\lambda=\{\alpha\in\Sigma\mid \alpha^\vee(\lambda)\in\Z\}$$
的 Weyl 群。置 $\Sigma_\lambda^+=\Sigma^+\cap\Sigma_\lambda$，就在 $\Sigma_\lambda$ 上定义了序。这给出 $\Sigma_\lambda$ 的一组单根\index{dangen@单根|hyperpage} $\Pi_\lambda$，以及相应的单反射\index{danfanshe@单反射|hyperpage}集 $S_\lambda$。记 $\ell_\lambda$ 为 $(W_\lambda,S_\lambda)$ 上的长度函数\index{changduhanshu@长度函数|hyperpage}。称 $\lambda\in\h^*$ 为\emph{正则}的，若对一切 $\alpha\in\Sigma$ 有 $\alpha^\vee(\lambda)\neq0$；称 $\lambda$ 为\emph{反支配\index{zhipei@支配|hyperpage}\index{fanzhipei@反支配|hyperpage}}的，若对一切 $\alpha\in\Sigma^+$，$\alpha^\vee(\lambda)$ 都不是严格正整数。置
$$n(\lambda)=\min\{\ell_\lambda(w)\mid w\in W_\lambda,\ w\lambda \text{ 反支配}\}.$$
特别地，$n(\lambda)=0$ 等价于 $\lambda$ 反支配。设 $P(\Sigma)$ 是 $\h^*$ 中的权格。

显然，$\mu\in P(\Sigma)$ 自然决定 $X$ 上的齐性可逆 $\calO_X$-模 $\calO(\mu)$。若 $V$ 是 $X$ 上的 $\D_\lambda$-模，则用可逆 $\calO_X$-模 $\calO(\mu)$ 去扭，得
$$V(\mu)=V\otimes_{\calO_X}\calO(\mu),$$
它是 $X$ 上的 $\D_{\lambda+\mu}$-模（……）。这个构造定义了从范畴 $\mathcal{M}(\D_\lambda)$ 到 $\mathcal{M}(\D_{\lambda+\mu})$ 的协变函子。称之为\emph{几何平移函子\index{pingyihanzi@平移函子|hyperpage}}（geometric translation functor）。它显然是范畴等价，并且也诱导 $\mathrm{M}_{qc}(\D_\lambda)$（相应地 $\mathrm{M}_{coh}(\D_\lambda)$）与 $\mathrm{M}_{qc}(\D_{\lambda+\mu})$（相应地 $\mathrm{M}_{coh}(\D_{\lambda+\mu})$）的等价。

几何平移与另一构造密切相关。设 $F$ 是有限维 $\g$-模。则层 $\mathcal{F}=\calO_X\otimes_\C F$ 带有自然的 $U^\circ$-模结构。下面定义它的一个滤过，它与模 $F$ 的权结构相关。

固定基点 $x_0\in X$。$\bb_{x_0}$-模 $F$ 有滤过
$$0=F_0\subset F_1\subset\cdots\subset F_m,\qquad m=\dim F,$$
使得
$$\dim(F_i/F_{i-1})=1,\qquad \n_{x_0}F_i\subset F_{i-1},\qquad 1\le i\le m.$$
因此 $\bb_{x_0}/\n_{x_0}$ 自然作用在 $F_i/F_{i-1}$ 上；经特殊化，这个作用诱导了 Cartan 代数 $\h$ 在 $F_i/F_{i-1}$ 上的作用，由某个权 $\nu_i\in P(\Sigma)$ 给出。

显然，$\nu_i<\nu_j$ 蕴含 $i>j$。层 $\mathcal{F}$ 是平凡齐性向量丛 $X\times F\to X$ 的局部截面层。$\mathcal{F}$ 的滤过诱导了这个向量丛的滤过，其子丛是齐性向量子丛，在基点 $x_0$ 处的纤维为 $F_i$（$1\le i\le m$）。记 $\mathcal{F}_i$（$1\le i\le m$）为这些子丛的局部截面层。它们是局部自由凝聚 $\calO_X$-模，也是 $U^\circ$-模。另一方面，作为 $U^\circ$-模有 $\mathcal{F}_i/\mathcal{F}_{i-1}=\calO(\nu_i)$，即 $\mathcal{F}_i/\mathcal{F}_{i-1}$ 自然是 $\D_{\nu_i-\rho}$-模。

设 $V$ 是 $X$ 上的拟凝聚\index{niningju@拟凝聚|hyperpage} $\D_\lambda$-模。则 $\calO_X$-模 $V\otimes_{\calO_X}\mathcal{F}$ 带有自然的 $U^\circ$-模结构，由
$$\xi(v\otimes s)=\xi v\otimes s+v\otimes\xi s,\qquad \xi\in\g$$
给出，其中 $v$ 与 $s$ 分别是 $V$ 与 $\mathcal{F}$ 的局部截面。用子模 $V\otimes_{\calO_X}\mathcal{F}_i$（$1\le i\le m$）可定义它的 $U^\circ$-模滤过。由上述讨论，相应的分次模是 $V(\nu_i)$（$1\le i\le m$）的直和。因此，对任意 $\xi\in\zz(\g)$，乘积 $\prod_{1\le i\le m}(\xi-\chi_{\lambda+\nu_i}(\xi))$ 零化 $V\otimes_{\calO_X}\mathcal{F}$。由初等线性代数，$V\otimes_{\calO_X}\mathcal{F}$ 分解为其 $\zz(\g)$-广义特征层的直和。

设 $V$ 是 $U^\circ$-模，$\lambda\in\h^*$。记 $V[\lambda]$ 为 $V$ 的、对应于 $\chi_\lambda$ 的 $\zz(\g)$-广义特征层。

\begin{lemma}\label{lem:02.2.2.1}
设 $\lambda\in\h^*$，$\mu\in P(\Sigma)$，$w\in W$，且 $w\lambda$ 与 $-w\mu$ 都反支配。设 $F$ 是以 $w\mu$ 为最高权的不可约有限维 $\g$-模。则
$$V\longmapsto\bigl(V(-\mu)\otimes_{\calO_X}\mathcal{F}\bigr)[\lambda]$$
是从 $\mathcal{M}(\D_\lambda)$ 到自身的协变函子，且自然等价于恒等函子。
\end{lemma}

\begin{proof}
$V(-\mu)\otimes_{\calO_X}\mathcal{F}$ 的滤过的合成因子是 $V(-\mu+\nu)$，其中 $\nu$ 跑遍 $F$ 的全部权。故 $\zz(\g)$ 以无穷小特征标 $\chi_{\lambda-\mu+\nu}$ 作用在它们上。设对某个 $s\in W$ 有
$$s\lambda=\lambda-\mu+\nu.$$
置 $s'=wsw^{-1}$，$\lambda'=w\lambda$，则
$$s'\lambda'-\lambda'=w\nu-w\mu.$$
又 $w\mu$ 与 $w\nu$ 都是 $F$ 的权，故 $s'\lambda'-\lambda'\in Q(\Sigma)$。因此 $s'\in W_{\lambda'}$。由于 $w\mu$ 是 $F$ 的最高权，$w\nu-w\mu$ 是若干负根之和。另一方面，由于 $\lambda'$ 反支配，$s'\lambda'-\lambda'$ 是 $\Sigma_\lambda^+\subset\Sigma^+$ 中若干根之和。于是 $s\lambda=\lambda$ 且 $\mu=\nu$；而 $V(-\mu)\otimes_{\calO_X}\mathcal{F}$ 中对应于 $\chi_\lambda$ 的广义特征层同构于 $V$。
\end{proof}

\begin{lemma}\label{lem:02.2.2.2}
设 $\lambda\in\h^*$，$\mu\in P(\Sigma)$，$w\in W$，且 $w\lambda$ 与 $-w\mu$ 都反支配。设 $\lambda$ 与 $\lambda-\mu$ 在 $W$ 中的稳定化子相同。设 $F$ 是以 $-w\mu$ 为最低权的不可约有限维 $\g$-模。则
$$V\longmapsto\bigl(V(\mu)\otimes_{\calO_X}\mathcal{F}\bigr)[\lambda-\mu]$$
是从 $\mathcal{M}(\D_{\lambda-\mu})$ 到自身的协变函子，且自然等价于恒等函子。
\end{lemma}

\begin{proof}
$V(\mu)\otimes_{\calO_X}\mathcal{F}$ 的滤过的合成因子是 $V(\mu+\nu)$，其中 $\nu$ 跑遍 $F$ 的全部权。故 $\zz(\g)$ 以无穷小特征标 $\chi_{\lambda+\nu}$ 作用在它们上。设对某个 $s\in W$ 有
$$\lambda-\mu=s(\lambda+\nu).$$
置 $s'=wsw^{-1}$，$\lambda'=w\lambda$，则
$$\lambda'-s'\lambda'=s'w\nu+w\mu.$$
由于 $-w\mu$ 是 $F$ 的最低权，$s'w\nu+w\mu$ 是若干正根之和。因此 $s'\in W_{\lambda'}$。又 $w\lambda$ 反支配，故 $s'\lambda'-\lambda'=0$，即 $s'\lambda'=\lambda'$；于是 $s\lambda=\lambda$。按假设，$s$ 稳定 $\lambda-\mu$，这意味着 $\nu=-\mu$。于是 $V(\mu)\otimes_{\calO_X}\mathcal{F}$ 中对应于 $\chi_{\lambda-\mu}$ 的广义特征层就是 $V$。
\end{proof}
\begin{annote}
底本在 2.2 的证明中把 $s'\in W_{\lambda'}$ 印作 $s\in W_{\lambda'}$。作用在 $\lambda$ 上的是 $s$，而 $W_{\lambda'}$ 的下标是 $\lambda'=w\lambda$；且下一句要由 $w\lambda$ 的反支配性推出 $s'\lambda'=\lambda'$。故按上下文应为 $s'$。[来源：底本 p.23，上下文自证]
\end{annote}

设 $\lambda\in\h^*$ 使 $n(\lambda)=k$，$k>0$。则存在 $w\in W_\lambda$，使 $\ell_\lambda(w)=k$ 且 $w\lambda$ 反支配。设
$$w=s_{\beta_1}s_{\beta_2}\cdots s_{\beta_k}$$
是 $w$ 在 $(W_\lambda,S_\lambda)$ 中的\emph{既约表达式}（reduced expression）。置 $\alpha=w^{-1}\beta_1$，$w'=s_{\beta_1}w$。则 $w'=s_{\beta_2}\cdots s_{\beta_k}$，且 $\ell_\lambda(w')=k-1$。于是
$$w's_\alpha\lambda=s_{\beta_1}ws_\alpha\lambda=s_{\beta_1}s_{w\alpha}w\lambda=w\lambda$$
反支配，故
$$n(s_\alpha\lambda)\le\ell_\lambda(w')=k-1.$$

现在，$w\lambda$ 的反支配性给出 $\beta_1^\vee(w\lambda)\in-\Z_+$；而 $\beta_1^\vee(w\lambda)=0$ 将推出 $w'\lambda=s_{\beta_1}w\lambda=w\lambda$ 反支配，与 $w$ 的选取矛盾。因此
$$p=-\beta_1^\vee(w\lambda)\in\N.$$
设 $C_\lambda$ 是 $\Sigma_\lambda^+$ 对应的 Weyl 房，则方程 $\beta_1^\vee(\tau)=0$ 确定 $C_\lambda$ 的一面墙。显然，$C_\lambda$ 的 $\Sigma$-正则点被分成有限多个 $\Sigma$ 的 Weyl 房，其中至少有一个与 $C_\lambda$ 共此墙。取 $C$ 为这样的一个 Weyl 房。

设 $\sigma\in P(\Sigma)\cap C$，使 $\beta_1^\vee(\sigma)=p$。则 $w\lambda-s_{\beta_1}\sigma$ 落在 $\beta_1^\vee$ 确定的墙中。又因
$$\Sigma_\lambda=\Sigma_{w\lambda}=\Sigma_{w\lambda-s_{\beta_1}\sigma}$$
以及
$$s_{\beta_1}\bigl(\Sigma_\lambda^+-\{\beta_1\}\bigr)=\Sigma_\lambda^+-\{\beta_1\},$$
可见对 $\beta\in\Sigma_\lambda^+-\{\beta_1\}$ 有
$$\beta^\vee(w\lambda-s_{\beta_1}\sigma)=\beta^\vee(w\lambda)-(s_{\beta_1}\beta)^\vee(\sigma)\in-\Z_+,$$
故 $w\lambda-s_{\beta_1}\sigma$ 反支配。于是，由
$$w'(\lambda-s_\alpha w^{-1}\sigma)=s_{\beta_1}w(\lambda-s_\alpha w^{-1}\sigma)=s_{\beta_1}(w\lambda-s_{\beta_1}\sigma)=w\lambda-s_{\beta_1}\sigma,$$
可得
$$n(\lambda-s_\alpha w^{-1}\sigma)\le\ell_\lambda(w')=k-1.$$

现在设 $V$ 是 $\D_\lambda$-模。则其平移 $V(p\alpha)$ 是 $\D_{\lambda+p\alpha}$-模。又有
$$\lambda+p\alpha=\lambda-\beta_1^\vee(w\lambda)\alpha=\lambda-\alpha^\vee(\lambda)\alpha=s_\alpha\lambda.$$
类似地，平移 $V(-s_\alpha w^{-1}\sigma)$ 是 $\D_{\lambda-s_\alpha w^{-1}\sigma}$-模。

设 $F$ 是以 $\sigma$ 为\emph{极端权}（extremal weight）的不可约有限维 $\g$-模。置
$$G=\bigl(V(-s_\alpha w^{-1}\sigma)\otimes_{\calO_X}\mathcal{F}\bigr)[\lambda].$$
则 $V(-s_\alpha w^{-1}\sigma)\otimes_{\calO_X}\mathcal{F}$ 的滤过诱导了 $G$ 的滤过
$$G_i=G\cap\bigl(V(-s_\alpha w^{-1}\sigma)\otimes_{\calO_X}\mathcal{F}_i\bigr),\qquad 1\le i\le m.$$
这个滤过具有下述性质：除「$\lambda$ 与 $\lambda-s_\alpha w^{-1}\sigma+\nu_i$ 落在同一个 Weyl 群轨道 $\theta$ 中」这一情形外，$G_i=G_{i-1}$。若该条件满足，则
$$G_i/G_{i-1}=V(-s_\alpha w^{-1}\sigma+\nu_i).$$
因此，为看清 $G$ 的结构，须找出 $F$ 的全部权 $\nu$，使
$$\lambda-s_\alpha w^{-1}\sigma+\nu=s\lambda$$
对某个 $s\in W$ 成立。这意味着
$$s\lambda-\lambda=\nu-s_\alpha w^{-1}\sigma\in Q(\Sigma),$$
故 $s\in W_\lambda$。于是 $s'=wsw^{-1}$ 满足
$$s'w\lambda-w\lambda=w(s\lambda-\lambda)=w(\nu-s_\alpha w^{-1}\sigma)=w\nu-s_{\beta_1}\sigma,$$
而 $w\lambda$ 反支配，故
$$s'(w\lambda)-w\lambda=\sum_{\beta\in\Pi_\lambda}m_\beta\beta,\qquad m_\beta\in\Z_+.$$
由此可见
$$s_{\beta_1}\sigma-w\nu=-\sum_{\beta\in\Pi_\lambda}m_\beta\beta.$$
按 $C$ 的选取，$\Sigma$ 中（关于 $C$ 所定义的序的）全部正根构成的集合 $\Sigma^+(C)$ 含 $\Sigma_\lambda^+$。因此，若记 $\Pi(C)$ 为 $C$ 决定的 $\Sigma$ 的单根集，则 $w\nu-s_{\beta_1}\sigma$ 是 $\Pi(C)$ 中若干根之和。由于根 $\beta_1$ 属于 $\Pi(C)$，
$$s_{\beta_1}w\nu-\sigma=s_{\beta_1}(w\nu-s_{\beta_1}\sigma)$$
是 $\Pi(C)$ 中若干根之和与 $r\beta_1$（$r\in\Z_+$）之差。另一方面，$s_{\beta_1}w\nu$ 是权，而 $\sigma$ 是 $F$ 关于 $C$ 所定义的序的最高权，故 $\sigma-s_{\beta_1}w\nu$ 是 $\Pi(C)$ 中若干根之和。最终得到：存在 $q\in\Z_+$，使
$$\sigma-s_{\beta_1}w\nu=q\beta_1.$$
\begin{annote}
底本此处把 $s_{\beta_1}(w\nu-s_{\beta_1}\sigma)$ 印作 $s_{\beta_1}(w\nu-s_{\beta_1})$，漏掉了 $\sigma$。由于 $s_{\beta_1}$ 是对合，下方要用的等式 $s_{\beta_1}w\nu-\sigma=s_{\beta_1}(w\nu-s_{\beta_1}\sigma)$ 只有补上 $\sigma$ 才成立。[来源：底本 p.23，上下文自证]
\end{annote}
于是
$$s'w\lambda-w\lambda=w\nu-s_{\beta_1}\sigma=-s_{\beta_1}(\sigma-s_{\beta_1}w\nu)=q\beta_1;$$
再引入 $\h^*$ 上标准的 $W$-不变双线性型，得
$$\|\lambda\|^2=\|s'w\lambda\|^2=\|w\lambda+q\beta_1\|^2=\|\lambda\|^2+2q(w\lambda|\beta_1)+q^2\|\beta_1\|^2.$$
这意味着 $q=0$ 或 $q=-\beta_1^\vee(w\lambda)=p$。

第一种情形下，$\nu=s_\alpha w^{-1}\sigma$ 是 $F$ 的极端权。于是当 $\nu_i=\nu$ 时 $G_i/G_{i-1}=V$，而这样的 $i$（$1\le i\le m$）只有一个。第二种情形下，
$$w\nu=s_{\beta_1}\sigma+p\beta_1=s_{\beta_1}(\sigma-p\beta_1)=\sigma,$$
故 $\nu=w^{-1}\sigma$ 仍是 $F$ 的极端权。此时
$$-s_\alpha w^{-1}\sigma+\nu=-s_\alpha w^{-1}\sigma+w^{-1}\sigma=\alpha^\vee(w^{-1}\sigma)\alpha=\beta_1^\vee(\sigma)\alpha=p\alpha,$$
这意味着当 $\nu_j=\nu$ 时 $G_j/G_{j-1}=V(p\alpha)$，而这样的 $j$（$1\le j\le m$）只有一个。因此 $U^\circ$-模 $G$ 有长度为 $2$ 的合成列\index{hechenglie@合成列|hyperpage}，其子商依次为 $V$ 与 $V(p\alpha)$。最后，由 [5, Ch. VI, Cor. 2 of Prop. 17]，
$$\alpha=w^{-1}\beta_1=s_{\beta_k}\cdots s_{\beta_1}\beta_1=-s_{\beta_k}\cdots s_{\beta_2}\beta_1\in-\Sigma_\lambda^+\subset-\Sigma^+,$$
由此得
$$\nu_j=w^{-1}\sigma<w^{-1}\sigma-p\alpha=w^{-1}\sigma-\beta_1^\vee(\sigma)\alpha=w^{-1}\sigma-\alpha^\vee(w^{-1}\sigma)\alpha=s_\alpha w^{-1}\sigma=\nu_i,$$
故由前述注记 $i<j$。这给出正合列
$$0\longrightarrow V\longrightarrow G\longrightarrow V(p\alpha)\longrightarrow 0.$$
